% !PRINTED 621
\weberpage{621}
% Faint lines above §181 are reversed verso show-through; no direct recto body text is present.

\begin{center}
\S.\,181.

\gesperrt{Application to the metacyclic equations}
\gesperrt{of degree} $5$.
\end{center}

\quad We make one further application of the general theorems
to the equations of degree $5$. If $n = 5$, the cyclic
group has $5$, the half-metacyclic $10$, and the full metacyclic
group $20$ permutations. For a general equation of degree $5$, a metacyclic
function satisfies a resolvent of degree $6$,
and if this equation of degree $6$ has a rational
root, then the given equation of degree $5$ is metacyclic.
The half-metacyclic functions, i.\,e. the functions
that admit the permutations of the half-metacyclic group,
satisfy a resolvent of degree $12$, which however, after
adjunction of the square root of the discriminant, splits into two
factors of degree $6$.

\quad If one applies the generating substitutions for the modulus $5$
\[
s = (z, z+1),\quad t = (z, 2z),\quad t^2 = (z, 4z)
\]
to the figures $0, 1, 2, 3, 4$, there result the permutations
\weq{1}{
\begin{array}{lccccc}
 & 0, & 1, & 2, & 3, & 4\\
(s) & 1, & 2, & 3, & 4, & 0\\
(t) & 0, & 2, & 4, & 1, & 3\\
(t^2) & 0, & 4, & 3, & 2, & 1.
\end{array}
}

\quad If we apply the substitutions $s, t^2, t$ to the pairs of
figures $(0,1), (1,2), (2,3), (3,4), (4,0)$, it follows that
\weq{2}{
\begin{array}{lccccc}
 & (0,1), & (1,2), & (2,3), & (3,4), & (4,0)\\
(s) & (1,2), & (2,3), & (3,4), & (4,0), & (0,1)\\
(t^2) & (0,4), & (4,3), & (3,2), & (2,1), & (1,0);
\end{array}
}
these pairs therefore remain unchanged, in their totality, by $s$ and
by $t^2$, and consequently by the half-metacyclic group,
and a symmetric function of the corresponding
pairs of roots is half-metacyclic. The substitution $t$ effects
the permutations
\weq{3}{
\begin{array}{lccccc}
 & (0,1), & (1,2), & (2,3), & (3,4), & (4,0)\\
(t) & (0,2), & (2,4), & (4,1), & (1,3), & (3,0),
\end{array}
}
%CONT>>
